% !TeX TS-program = pdflatex

\documentclass[a4paper]{article}

% \usepackage[default]{fontsetup}

\usepackage{fancyhdr}
\usepackage{extramarks}
\usepackage{amsmath}
\usepackage{amsthm}
\usepackage{amsfonts}
\usepackage{tikz}
\usepackage[plain]{algorithm}
\usepackage{algpseudocode}
\usepackage{enumerate}

\usetikzlibrary{automata,positioning}

%
% Basic Document Settings
%  

\topmargin=-0.2in
\evensidemargin=0in
\oddsidemargin=0in
\textwidth=6.5in
\textheight=9.5in
\headsep=0.25in

\linespread{1.1}

\pagestyle{fancy}
\lhead{\hmwkAuthorName}
\chead{\hmwkClass : \hmwkTitle}
\rhead{\firstxmark}
\lfoot{\lastxmark}
\cfoot{\thepage}

\renewcommand\headrulewidth{0.4pt}
\renewcommand\footrulewidth{0.4pt}

\setlength\parindent{0pt}

%
% Create Problem Sections
%

%修改下标的地方
%\newcommand{\enterProblemHeader}[1]{
	%	\nobreak\extramarks{}{Problem \arabic{#1} continued on next page\ldots}\nobreak{}
	%	\nobreak\extramarks{Problem \arabic{#1} (continued)}{Problem \arabic{#1} continued on next page\ldots}\nobreak{}
	%}

\newcommand{\exitProblemHeader}[1]{
    \nobreak\extramarks{Problem \arabic{#1} }{Problem \arabic{#1} continued on next page\ldots}\nobreak{}
	\stepcounter{#1}
	\nobreak\extramarks{Problem \arabic{#1}}{}\nobreak{}
}

\newcommand*\circled[1]{\tikz[baseline=(char.base)]{
		\node[shape=circle,draw,inner sep=2pt] (char) {#1};}}


\setcounter{secnumdepth}{0}
\newcounter{partCounter}
\newcounter{homeworkProblemCounter}
\setcounter{homeworkProblemCounter}{1}
\nobreak\extramarks{Problem \arabic{homeworkProblemCounter}}{}\nobreak{}

%
% Homework Problem Environment
%
% This environment takes an optional argument. When given, it will adjust the
% problem counter. This is useful for when the problems given for your
% assignment aren't sequential. See the last 3 problems of this template for an
% example.
%

\newenvironment{homeworkProblem}[1][-1]{
	\ifnum#1>0
	\setcounter{homeworkProblemCounter}{#1}
	\fi
	\section{Problem \arabic{homeworkProblemCounter} (25\%)}
	\setcounter{partCounter}{1}
}{
	\exitProblemHeader{homeworkProblemCounter}
}

%
% Homework Details
%   - Title
%   - Class
%   - Due date
%   - Name
%   - Student ID

\newcommand{\hmwkTitle}{Homework\ \#02}
\newcommand{\hmwkClass}{Probability \& Statistics for EECS}
\newcommand{\hmwkDueDate}{2026-10-4}
\newcommand{\hmwkAuthorName}{}
\newcommand{\hmwkAuthorID}{}


%
% Title Page
%

\title{
	\vspace{2in}
	\textmd{\textbf{\hmwkClass:\\  \hmwkTitle}}\\
	\normalsize\vspace{0.1in}\small{Due\ on\ \hmwkDueDate\ at 23:59}\\
	\vspace{4in}
}

\author{
	Name: \textbf{\hmwkAuthorName} \\
	Student ID: \hmwkAuthorID}
\date{}

\renewcommand{\part}[1]{\textbf{\large Part \Alph{partCounter}}\stepcounter{partCounter}\\}

%
% Various Helper Commands
%

% Useful for algorithms
\newcommand{\alg}[1]{\textsc{\bfseries \footnotesize #1}}
% For derivatives
\newcommand{\deriv}[1]{\frac{\mathrm{d}}{\mathrm{d}x} (#1)}
% For partial derivatives
\newcommand{\pderiv}[2]{\frac{\partial}{\partial #1} (#2)}
% Integral dx
\newcommand{\dx}{\mathrm{d}x}
% Alias for the Solution section header
\newcommand{\solution}{\textbf{\large Solution}}
% Probability commands: Expectation, Variance, Covariance, Bias
\newcommand{\E}{\mathrm{E}}
\newcommand{\Var}{\mathrm{Var}}
\newcommand{\Cov}{\mathrm{Cov}}
\newcommand{\Bias}{\mathrm{Bias}}

\begin{document}
	
	
	% \maketitle
	% \thispagestyle{empty}
	% \pagebreak
	
	\date{
		Due on Oct. 4, 2026, 23:59 UTC+8}
	\title{SI 140A-02  Probability \& Statistics for EECS, Fall 2026 \\
		Homework 2}
	\maketitle
	Read all the instructions below carefully before you start working on the assignment, and before you make a submission.
	\begin{itemize}
		\item You are required to write down all the major steps towards making your conclusions; otherwise you may obtain limited points of the problem.
		\item Write your homework in English; otherwise you will get no points of this homework.
		\item Any form of plagiarism will lead to $0$ point of this homework. 
        \item You are required to \textbf{submit your homework in PDF format generated by LaTeX}, otherwise you will
\textbf{get no points of this homework}.
	\end{itemize}

	\newpage
		
\begin{homeworkProblem}

A box contains 5 red balls, 4 blue balls, and 3 green balls. Three balls are drawn sequentially without replacement.

Let
$
A=\{\text{the first ball is red}\},
$
and let
$
C=\{\text{at least one of the three drawn balls is red}\}.
$

Suppose that the only information revealed after the experiment is that event $C$ occurred; that is, at least one of the three drawn balls was red.

\begin{enumerate}
    \item[(a)] Find $P(C)$.

    \item[(b)] Find $P(A\cap C)$ and hence calculate $P(A\mid C)$.

    \item[(c)] Let $B$ be the event that the second ball is red. Calculate $P(B\mid C)$.

    \item[(d)] Calculate $P(A\cap B\mid C)$.

    \item[(e)] Compare $P(A\mid C)$ with $P(A)$. Are $A$ and $C$ independent? Justify your answer using the definition of independence.
\end{enumerate}

\end{homeworkProblem}

\newpage

\begin{homeworkProblem}
\textbf{Multiple Detection System}

A monitoring system uses four independent sensors to detect whether a machine has experienced a failure. Let $A_i$ denote the event that sensor $i$ raises a false alarm.

The false-alarm probabilities are
$
P(A_1)=0.10,\;
P(A_2)=0.15,\;
P(A_3)=0.20,\;
P(A_4)=0.25.
$

Assume that the four events are mutually independent.

The monitoring system issues a warning if \textbf{at least one sensor} raises a false alarm.

\begin{enumerate}[(a)]
    \item Suppose the monitoring system has issued a warning. Find the probability that \textbf{at least two sensors} have raised a false alarm, i.e.,
    $
    P\left(
    \text{at least two alarms}
    \,\middle|\,
    A_1\cup A_2\cup A_3\cup A_4
    \right).
    $

    \item Suppose now that no independence assumption is made. Instead, you are given
    $
    P(A_i)=0.1,\; i=1,2,3,4,
    $
    and
    $
    P(A_i\cap A_j)=0.02
    $
    for every $i\neq j$, together with
    $
    P(A_i\cap A_j\cap A_k)=0.005
    $
    for every distinct $i,j,k$, while
    $
    P(A_1\cap A_2\cap A_3\cap A_4)=0.001.
    $

    Use the inclusion-exclusion principle to calculate
    $
    P(A_1\cup A_2\cup A_3\cup A_4).
    $
\end{enumerate}

\end{homeworkProblem}

\newpage

\begin{homeworkProblem}

A factory has three production lines, $L_1$, $L_2$, and $L_3$, which produce $30\%$, $45\%$, and $25\%$ of the factory's total products, respectively. The probabilities that a product is defective given that it comes from $L_1$, $L_2$, or $L_3$ are
$
P(D\mid L_1)=0.02,\;
P(D\mid L_2)=0.03,\;
P(D\mid L_3)=0.05,
$
where $D$ denotes the event that a randomly selected product is defective.

\begin{enumerate}

    \item[(a)] Two products are selected independently from the factory. Find the probability that both products are defective.

    \item[(b)] Three products are selected independently from the factory. Find the probability that at least one of the three products is defective.

    \item[(c)] Suppose $n$ products are selected independently from the factory. Find the smallest integer $n$ such that the probability of obtaining at least one defective product is at least $95\%$.

\end{enumerate}

\end{homeworkProblem}

\newpage

\begin{homeworkProblem}

Two dice are rolled independently. For each die, the probability of
showing an even number is $p$, where $0<p<1$. The two dice have the
same distribution, but no assumption is made that the individual
outcomes $1,2,\ldots,6$ are equally likely.

Define
$
A=\{\text{the sum of the two dice is even}\},
$
and
$
B=\{\text{the sum of the two dice is greater than }8\}.
$

Determine whether there exists a value of $p$ such that, for every
possible distribution of the individual die outcomes satisfying
$P(\text{even})=p$, the events $A$ and $B$ are independent.

If such a value of $p$ exists, find all such values and prove that the
independence holds for every admissible distribution.

If no such value exists, prove this by constructing, for each candidate
value of $p$, two admissible distributions that have the same value of
$p$ but give different values of
$
P(A\cap B)-P(A)P(B).
$

Your argument should distinguish carefully between independence that is
implied by the given information and independence that occurs only for
a particular choice of the underlying distribution.

\end{homeworkProblem}
	
	
\end{document}